% main.tex — نسخه‌ی نهایی
% ============================================================

\input{preamble}

\begin{document}

% ---- Title Page ----
\input{titlepage}

% ---- Dedication ----
% ═══════════════════════════════════════════════════════════════════════════
%  تقدیم
% ═══════════════════════════════════════════════════════════════════════════
\thispagestyle{empty}
\vspace*{6cm}
\begin{center}
  {\Large\itshape\color{PrimaryMid}
  تقدیم به همه‌ی ملت‌هایی که\\[8pt]
  در جست‌وجوی آزادی\\[8pt]
  بهای گزافی پرداخته‌اند.}
\end{center}
\vspace*{\fill}
\cleardoublepage

% ---- Front Matter ----
\frontmatter
\tableofcontents
\clearpage
\listoffigures
\clearpage
\listoftables

% ---- Preface ----
\chapter*{پیش‌گفتار}
\addcontentsline{toc}{chapter}{پیش‌گفتار}

\begin{quotebox}
«هنر حکمرانی، هنر زیستن با تفاوت‌هاست.»
\hfill — ویل کیملیکا
\end{quotebox}

\vspace{6pt}

این کتاب تلاشی است برای پاسخ به پرسشی بنیادین در سیاست تأسیسی:
چگونه می‌توان ساختاری سیاسی طراحی کرد که تنوع قومی و ملّی را نه تهدید، بلکه فرصت بشمارد؟

در بخش نخست، مبانی نظری شکل‌گیری دولت و رابطه‌ی آن با مفاهیم ملت و قومیت واکاوی می‌شود.
بخش دوم، الگوهای سیاسی مدیریت تنوع — از تمرکزگرایی مطلق تا کنفدراسیون — را با نمونه‌های
عینی بررسی می‌کند. بخش سوم مطالعات تطبیقی مفصلی از هندوستان، بریتانیا، سوئیس، بلژیک
و چند کشور دیگر ارائه می‌دهد. و سرانجام بخش چهارم، به‌طور ویژه به ایران اختصاص یافته
و شامل تحلیل وضعیت، چالش‌ها و نقشه‌ی راه اجرایی است.

این اثر نه مدعی ارائه‌ی نسخه‌ی واحد است و نه قصد ایدئولوژی‌سازی دارد. هدف آن ارائه‌ی
ابزارهای تحلیلی و تجربه‌های جهانی برای غنی‌ترکردن گفت‌وگوی ملّی درباره‌ی آینده‌ی ایران است.

\simplesep

\begin{keybox}[ساختار کتاب]
\begin{description}[style=nextline, leftmargin=2cm]
  \item[\textcolor{PrimaryDark}{بخش اول: مبانی نظری و مفهومی}]
    فصل ۱ (مبانی شکل‌گیری دولت) + فصل ۲ (الگوهای رابطه‌ی دولت–ملت–قومیت)
  \item[\textcolor{PrimaryDark}{بخش دوم: الگوها و مطالعات تطبیقی}]
    فصل ۳ (طیف رژیم‌های سیاسی) + فصل ۴ (مطالعات موردی: هند، بریتانیا و...)
  \item[\textcolor{PrimaryDark}{بخش سوم: ایران — تحلیل و راهبرد}]
    فصل ۵ (تحلیل ساختار قومی–ملّی ایران) + فصل ۶ (نقشه‌ی راه اجرایی)
  \item[\textcolor{PrimaryDark}{پیوست‌ها}]
    واژه‌نامه، نقشه‌ی مفهومی، منابع، جدول جامع مقایسه‌ای، چک‌لیست اجرایی
\end{description}
\end{keybox}

\vspace{1cm}
\begin{flushleft}
\textcolor{PrimaryDark}{\bfseries پژوهشگر}\\
بهار ۱۴۰۴
\end{flushleft}

% ===========================================================
\mainmatter
% ===========================================================

% ---- PART I ----
\partpage{بخش اول}{مبانی نظری و مفهومی}

\input{chapter1}
\input{chapter2}

% ---- PART II ----
\partpage{بخش دوم}{الگوها و مطالعات تطبیقی}

\input{chapter3}
\input{chapter4}

% ---- PART III ----
\partpage{بخش سوم}{ایران: تحلیل و راهبرد}

\input{chapter5}
\input{chapter6}

% ---- Back Matter ----
\backmatter

\input{appendix}

\printbibliography[
  heading=bibintoc,
  title={کتاب‌نامه}
]

\printindex

% ---- Colophon ----
\input{colophon}

\end{document}